\documentclass[journal]{IEEEtran}
\usepackage{amsmath,amsfonts}
\usepackage{mathtools}
\usepackage{hyperref}
\usepackage{array}
\usepackage[caption=false,font=normalsize,labelfont=sf,textfont=sf]{subfig}
\usepackage{textcomp}
\usepackage{stfloats}
\usepackage{xurl}
\usepackage{graphicx}
\usepackage{siunitx}
\usepackage{cite}
\usepackage{orcidlink}
\usepackage{svg}
\usepackage{graphicx}

\hyphenation{op-tical net-works semi-conduc-tor IEEE-Xplore}

\begin{document}

\title{Task-Based Contingency Assessment of Power Networks: Numerical Coverage and Parallel Performance}

\author{Alexandre~Gràcia-Calvo\orcidlink{0009-0009-6728-6713},
        
        Eduardo~Iraola\orcidlink{0000-0002-0837-7086}, 
        Carlos~Herrera~Vázquez\orcidlink{0009-0002-2181-9683},
        Francesca~Rossi\orcidlink{0000-0002-4791-151X},
        Juan~Carlos~Olives-Camps\orcidlink{0000-0002-3122-1223}, \IEEEmembership{Member, IEEE,}
        and~Eduardo~Prieto-Araujo\orcidlink{0000-0003-4349-5923}, \IEEEmembership{Senior~Member, IEEE}% <-this % stops a space
\thanks{The authors are with CITCEA-UPC and BSC, Barcelona, Spain. E-mail: alexandre.gracia@upc.edu, carlos.herrera.vazquez@upc.edu, eduardo.iraola@bsc.edu, francesca.rossi@upc.edu, juan.carlos.olives@upc.edu, eduardo.prieto-araujo@upc.edu}% <-this % stops a space
\thanks{Manuscript received November X, 2025; revised December X, 2025.}}

% Headers
\markboth{IEEE Transactions on Power Systems,~Vol.~X, No.~X, November~2025}%
{Gr\`acia-Calvo \MakeLowercase{\textit{et al.}}: Critical Components in the IEEE 118-Bus Network}

\maketitle


\begin{abstract}
Exhaustive contingency assessment combines a rapidly growing set of outage scenarios with numerical analyses whose cost and applicability vary across cases. This paper investigates a task-based workflow for joint topological, power-flow and small-signal assessment of simultaneous single and double outages. The implementation uses PyCOMPSs to distribute independent contingency evaluations while a single master manages transactional persistence and subsequent risk aggregation. Each evaluation retains the supplied network configuration except for the disconnected elements, and distinguishes observed security failures from unresolved numerical or modelling outcomes. The experimental programme separates numerical verification, application coverage and parallel performance, with the IEEE 118-bus network as the main case study, containing 28,680 unique single and double outages. Performance evaluation will quantify end-to-end execution time, task-cost variability, resource efficiency and persistence overhead. Scalability and complete modal coverage remain experimental questions; neither is inferred from contingency counts alone.
\end{abstract}

\begin{IEEEkeywords}
Contingency analysis, high-performance computing, task-based workflows, PyCOMPSs, small-signal stability, risk assessment.
\end{IEEEkeywords}

\section{Introduction}
\IEEEPARstart{E}{xhaustive} contingency assessment requires evaluating combinations of component outages without conflating the security of the network with the numerical tools' ability to assess it. For a network with $m$ eligible elements, single and simultaneous double outages produce $m+\binom{m}{2}$ distinct cases. Even at a fixed contingency order, their computational costs can differ substantially: some cases fail a preliminary check, whereas others require electromagnetic transient (EMT) initialization and a modal calculation.

Parallel contingency assessment and high-performance dynamic security assessment are established research areas. Existing work includes massively parallel dynamic security platforms~\cite{konstantelos2017} and frameworks connecting simulation, uncertainty analysis and visualization~\cite{chen2018framework}. Consequently, distributing independent outages or combining severity with component reliability is insufficient, by itself, to establish a new contribution.

The focus here is the implementation and experimental assessment of a task-based workflow that combines heterogeneous numerical paths with centralized result management. PyCOMPSs~\cite{pycompss_softx,pycompss_tejedor2017} provides the distributed execution mechanism. The application preserves the supplied operating configuration, records why a case cannot reach a numerical stage, and propagates unresolved classifications into the aggregate assessment. This connects the computational measurements to the scientific coverage actually achieved.

The study addresses three questions: how variable numerical paths affect task costs and parallel efficiency; how result collection and transactional persistence affect end-to-end runtime and restart behaviour; and how incomplete numerical coverage limits the interpretation of component vulnerability. These are research questions to be answered by experiments, rather than claims of demonstrated novelty or scalability.

The manuscript combines the power-system use case and the computational study in a single contribution. The intended evidence comprises numerical verification, comparisons with a sequential execution of the same calculation, a controlled scaling study, and an explicit accounting of unresolved cases. No real-time operating capability is claimed before these measurements are available.

\section{Contingency Assessment and Scientific Semantics}
\label{sec:method}
\subsection{Unique Simultaneous Outages}
The eligible elements are AC lines, two-winding transformers and generators. Each case is generated from the original network by disabling its selected elements. Unordered pairs $\{i,j\}$ are evaluated once; they describe simultaneous outages, not a time-ordered cascading sequence. The remaining device parameters, slack assignments and operating setpoints are retained. No optimal-power-flow redispatch or artificial network repair is applied.

The current 118-bus input contains 175 lines, 11 two-winding transformers and 53 generators. Thus $m=239$, giving 239 single outages and 28,441 double outages, for a total of 28,680 cases. The base case is a separate validation experiment and is excluded from this catalogue.

\subsection{Numerical Path and Outcome Classification}
An evaluation records connectivity and balanced AC power-flow results. A connection-count guard determines whether the supplied network meets the EMT model builder's requirements. This guard is a model-applicability check: rejection does not, by itself, demonstrate electrical islanding or instability. Cases reaching the remaining stages undergo three-phase power flow, EMT initialization and convergence checks, followed by a Floquet-based modal calculation.

The adopted security criteria classify multiple electrical islands or an observed nonconverged power flow as a failure. Power-flow nonconvergence is an operational criterion of this study, not a mathematical proof that no feasible solution exists. A numerical exception without an observed convergence result is unresolved. When modal results are available, an exponent with real part above $10^{-6}\ \mathrm{s}^{-1}$ is classified as modal instability. Passing the computed-mode criterion together with the EMT health checks gives the workflow's stable classification; it does not certify the entire spectrum or transient stability.

The calculation requests 12 modes through the standard ARPACK path of the Floquet solver. Voltage and thermal limit violations are not included in the current classification. Generic dynamic models attached to a static network must be distinguished from an independently validated dynamic benchmark.

\subsection{Coverage-Aware Component Vulnerability}
Let $S_i$ and $S_{ij}$ denote binary severities for assessed single and double outages. The current component summaries are
\begin{equation}
R_i^{(1)}=S_i, \qquad
R_i^{(2)}=\sum_{j\ne i}p_j S_{ij},
\end{equation}
where $p_j$ is an assumed component-type weight. This second quantity describes weighted vulnerability conditional on the outage of $i$. It is not an annual system-failure frequency. The current generic weights are illustrative assumptions, rather than asset-specific reliability measurements.

For unresolved severities, lower and upper bounds are obtained by assigning zero and one, respectively, while retaining measured outcomes. Reports include both case coverage and weighted coverage. These bounds describe incomplete assessment; they are not statistical confidence intervals. A component ranking must be interpreted together with its unresolved contribution.

\section{Task-Based Execution and Result Management}
\subsection{Current Task Granularity}
The master loads contingency definitions and submits one PyCOMPSs task per case. A task loads the network, applies the outage, executes the applicable numerical stages and returns a structured result. All numerical stages currently execute inside that task. The implementation therefore exposes parallelism across contingencies; it does not yet expose a separate task graph for topology, power flow and modal analysis.

The independent cases form a naturally parallel workload. Its potential scheduling difficulty comes from variation in numerical paths and task durations, which must be measured. A task's requested computing units reserve runtime resources; they do not establish that its numerical kernels efficiently use that number of cores. Resource settings therefore require controlled comparison.

\subsection{Persistence and Resumption}
Workers do not access SQLite. The master receives results and writes transactions in batches, avoiding concurrent worker writes to the shared database. The batch-size setting controls results per transaction, not the number of concurrently scheduled tasks. Tasks are submitted before the master collects successive batches in submission order. Consequently, a slow case can delay persistence of its batch even when later results are ready; the performance study should quantify this effect.

Committed batches survive interruption, and a subsequent pending-only run selects unfinished cases. This provides application-level resumption at transaction boundaries. It is distinct from transparent runtime recovery of an executing task. Futures are released after their results have been persisted.

\subsection{Candidate Workflow Extensions}
Separating inexpensive checks from EMT and modal tasks could expose stage dependencies and allow stage-specific resource requirements. Any such extension must be compared with the existing one-task-per-case baseline using identical physical models and outcome rules. Its usefulness depends on scheduling, serialization and model-transfer costs; more tasks do not automatically imply better performance. This extension is proposed experimental work and is not part of the currently running implementation.

\section{Experimental Design}
\subsection{Numerical Verification and Application Coverage}
The IEEE 118-bus network is the main case study for contingency outcomes and parallel performance. Small-system checks support numerical verification: an analytical stable/unstable scalar reference and a 9-bus network used for local debugging, time-step refinement, periodic-trajectory closure, eigenpair residuals and comparison with a dense eigensolver applied to the same discrete operator. These checks do not replace verification of initialization, periodicity and modal accuracy on the 118-bus network, nor do they independently validate its physical dynamic model.

The IEEE 118-bus case assesses unique-contingency enumeration and application coverage at a larger catalogue size. Preliminary checks identify a bus, SARGENTS, that does not meet the EMT connection-count requirement in the supplied base network. This limitation must be reported without modifying the network to force acceptance. A workload that stops before EMT cannot be used as evidence for the scaling of the full modal workflow. The final report must show how many cases reach every numerical stage.

\subsection{Performance Protocol}
The same fixed case catalogue, model and numerical tolerances must be used when comparing resource configurations. Report end-to-end application time, calculation-stage time, task-duration distributions, completed cases per second, master persistence time and allocated node-hours. Scheduler queue waiting time is reported separately from execution time. Distinguish cases completing their assigned evaluation from those obtaining a resolved security classification.

For a fixed workload and a measured one-node baseline, define
\begin{equation}
S(n)=\frac{T(1)}{T(n)}, \qquad E(n)=\frac{S(n)}{n}.
\end{equation}
The same per-node settings are required for this comparison. A sequential execution of the same calculations provides a separate algorithmic baseline where practical. Compare task resource requests and persistence batch sizes independently. Repeat measurements sufficiently to report variability and retain the distribution of numerical outcomes so that runtime differences cannot be mistaken for equivalent scientific coverage.

The planned 20-node application run was cancelled before execution to prioritize modal validation; it provides no performance measurement. A subsequent study should begin with a small geometric sequence of node counts and extend only while there is enough work per worker. If the 118-bus catalogue cannot exercise the modal stages, select an additional documented dynamic case. Increasing contingency order alone cannot repair model incompatibility or validate full-workflow scaling.

\section{IEEE 118-Bus Results and Performance Evaluation}
\subsection{Supporting Numerical Verification}
An analytical scalar reference and a 9-bus network were used for development and numerical checks. The scalar reference recovers stable and unstable signs with approximately first-order exponent convergence under the current tangent-history initialization. The 9-bus checks show sensitivity of periodic closure to time step and simulated duration, and compare selected eigenpairs with the dense spectrum of the same discrete operator. Detailed development measurements are retained in the project documentation. They do not establish numerical accuracy or physical validity for the IEEE 118-bus case.

\subsection{118-Bus Base-Case Verification}
This subsection remains pending resolution of the EMT connection-validation issue. Before contingency results are interpreted, the 118-bus base case must demonstrate EMT initialization and convergence, adequate periodic closure, and modal checks under time-step and simulated-duration refinement. The periodicity criterion and interpretation of nearly neutral modes must be justified for this model.

\subsection{Contingency Outcomes and Coverage}
Results for the 239 single and 28,441 unique double outages remain pending. The final analysis will report coverage by numerical stage, observed security failures, unresolved cases and component vulnerability with its coverage bounds. Small-system development checks are excluded from these application results.

\subsection{Parallel Performance}
Controlled measurements on the 118-bus workload remain pending. They will compare end-to-end runtime, task-cost distributions, resource efficiency and persistence overhead for the same fixed catalogue and numerical settings. A workload blocked before EMT cannot establish modal-workflow performance. If the 118-bus case does not provide enough work for larger allocations, the reported scaling range will be limited accordingly.

\section{Discussion and Limitations}
A combined power-systems and HPC paper needs both scientifically interpretable outputs and computational evidence beyond the use of a parallel runtime. The present workflow provides an experimental baseline; whether finer task granularity or resource-aware execution improves it remains to be established. The original operating configuration and numerical acceptance rules must remain fixed across these comparisons.

Application coverage, generic dynamic-model assumptions, selected-mode analysis and illustrative component weights limit the conclusions that can be drawn. These limitations also determine which performance experiments represent the intended application. Large contingency counts alone establish neither a large dynamic workload nor a novel HPC method.

\section{Conclusion}
\label{conclusion}
The manuscript is being developed around a unified application and workflow study: exhaustive contingency assessment with explicit numerical coverage, task-based execution and resumable result management. Final conclusions on scalability, resource efficiency and component vulnerability are deferred until the numerical validation and controlled performance experiments are complete.

\section*{Acknowledgment}
\label{acknowledgement}

The Project TED2021-130351B-C21 (HP2C-DT) is funded by MICIU/AEI/10.13039/501100011033 and by the European Union NextGenerationEU/PRTR.


\bibliographystyle{IEEEtran}
\begin{thebibliography}{99}

\bibitem{chen2018framework}
Y.~Chen, B.~J.~Palmer, P.~Sharma, Y.~Yuan, B.~Mathew, and Z.~Huang,
``A High Performance Computational Framework for Dynamic Security Assessment under Uncertainty,''
in \emph{IEEE Electronic Power Grid (eGrid)}, 2018.
doi: \href{https://doi.org/10.1109/eGRID.2018.8598684}{10.1109/eGRID.2018.8598684}.



\bibitem{Breton2025}
J.~Breton, J.~Jaskolka, and G.~O.~M.~Yee,
``An approach to determine a system's behavioral security posture''
in \emph{Foundations and Practice of Security},
K.~Adi, S.~Bourdeau, C.~Durand, V.~Viet Triem Tong, A.~Dulipovici, 
Y.~Kermarrec, and J.~Garcia-Alfaro, Eds.
Cham, Switzerland: Springer Nature Switzerland, 2025, pp.~94--110.
doi: \href{https://doi.org/10.1007/978-3-031-87499-4_7}{10.1007/978-3-031-87499-4\_7}.

\bibitem{entsoe_28apr2025}
European Network of Transmission System Operators for Electricity (ENTSO-E),
``Grid Incident in Spain and Portugal on 28 April 2025,''
Public report, Oct. 2025.
[Online]. Available:
\url{https://eepublicdownloads.blob.core.windows.net/public-cdn-container/clean-documents/Publications/2025/entso-e_incident_report_ES-PT_April_2025_06.pdf}

\bibitem{stability_2023}
Y. Gu and T. C. Green, 
``Power System Stability With a High Penetration of Inverter-Based Resources,'' 
in \emph{Proceedings of the IEEE}, vol.~111, no.~7, pp.~832-853, July 2023, 
doi: \href{https://doi.org/10.1109/JPROC.2022.3179826}{10.1109/JPROC.2022.3179826}.

\bibitem{billinton1994reliability}
R.~Billinton and R.~N.~Allan, \emph{Reliability Evaluation of Power Systems}, New York: Springer, 1996,
doi: \href{https://doi.org/10.1007/978-1-4899-1860-4}{10.1007/978-1-4899-1860-4}.

\bibitem{ieee_rts96}
Reliability Test System Task Force of the Application of Probability Methods Subcommittee, ``The IEEE Reliability Test System-1996,'' \emph{IEEE Transactions on Power Systems}, vol. 14, no. 3, pp. 1010--1020, Aug. 1999,
doi: \href{https://doi.org/10.1109/59.780914}{10.1109/59.780914}.

\bibitem{agreira}
C. I. F. Agreira, S. M. F. de Jesus, S. L. de Figueiredo, C. M. Ferreira, J. A. D. Pinto i F. P. M. Barbosa, 
“Probabilistic Steady-State Security Assessment of an Electric Power System Using a Monte Carlo Approach,” 
en \emph{Proceedings of the 41st International Universities Power Engineering Conference}, 2006, vol. 2, pp. 408-411. 
doi: \href{https://doi.org/10.1109/UPEC.2006.367509}{10.1109/UPEC.2006.367509}.

\bibitem{perkin}
S.~Perkin, C.~Hamon, R.~Kristjansson, H.~Stefansson and P.~Jensson,
``Framework for trajectory-based probabilistic security assessment of power systems,''
\emph{IET Generation, Transmission \& Distribution}, 
vol.~13, no.~7, pp.~1088--1094, Mar. 2019,
doi: \href{https://doi.org/10.1049/iet-gtd.2018.5396}{10.1049/iet-gtd.2018.5396}.

\bibitem{PSA_1_Pereira}
E.~J.~da~S.~Pereira, J.~T.~Pinho, M.~A.~B.~Galhardo, and W.~N.~Macêdo,
``Methodology of risk analysis by Monte Carlo Method applied to power generation with renewable energy,''
\emph{Renewable Energy}, vol.~69, pp.~347--355, Sep. 2014.
doi: \href{http://dx.doi.org/10.1016/j.renene.2014.03.054}{10.1016/j.renene.2014.03.054}.

\bibitem{shahzad}
U.~Shahzad and S.~Asgarpoor,
``Probabilistic risk assessment of an active distribution network using Monte Carlo simulation approach,''
in \emph{Proc. North American Power Symposium (NAPS)}, 2019, pp.~1--6.
doi: \href{https://doi.org/10.1109/NAPS46351.2019.9000225}{10.1109/NAPS46351.2019.9000225}.

\bibitem{donde}
V.~Donde, V.~Lopez, B.~C.~Lesieutre, A.~Pinar, C.~Yang and J.~Meza,
``Severe multiple contingency screening in electric power systems,''
\emph{IEEE Transactions on Power Systems}, 
vol.~23, no.~2, pp.~406--417, May 2008,
doi: \href{https://doi.org/10.1109/TPWRS.2008.919243}{10.1109/TPWRS.2008.919243}.

\bibitem{contingency_screening_islanding}
X. Zheng, S. Kotamarty, K. Thuerwaechter, A. Lee, S. H. Huang, and L. Xie,
``A Generalized Approach to Contingency Screening with System Islanding,''
doi: \href{https://hdl.handle.net/10125/102973}{10.24251/HICSS.2023.342}.

\bibitem{bu}
S.~Q.~Bu, W.~Du and H.~F.~Wang,
``Probabilistic analysis of small-signal stability of power systems, a survey,''
in \emph{Proc. Int. Conf. on Sustainable Power Generation and Supply (SUPERGEN 2012)}, 
Hangzhou, China, Jan. 2012,
doi: \href{https://www.researchgate.net/publication/261020926_Probabilistic_analysis_of_small-signal_stability_of_power_systems_-_A_survey}{10.1049/cp.2012.1830}.

\bibitem{tabrizchi}
A.~M.~Tabrizchi and M.~M.~Rezaei,
``Probabilistic small-signal stability analysis of power systems based on Hermite polynomial approximation,''
\emph{SN Applied Sciences}, 
vol.~3, no.~9, art.~784, 2021,
doi: \href{https://doi.org/10.1007/s42452-021-04765-4}{10.1007/s42452-021-04765-4}.

\bibitem{arevalo_sstoolbox}
J.~Arévalo-Soler, D.~Moutevelis, E.~Mateu-Barriendos, O.~Alican,
C.~Collados-Rodr\'iguez, M.~Cheah-Ma\~ne, E.~Prieto-Araujo, and O.~Gomis-Bellmunt,
``A Matlab-based Toolbox for Automatic EMT Modeling and Small-Signal Stability Analysis of Modern Power Systems,''
arXiv:2506.22201, 2025.
[Online]. Available: \url{https://arxiv.org/abs/2506.22201}

\bibitem{konstantelos2017}
I.~Konstantelos, G.~Jamgotchian, S.~H.~Tindemans, P.~Duchesne, S.~Cole, C.~Merckx, G.~Strbac, and P.~Panciatici,
``Implementation of a Massively Parallel Dynamic Security Assessment Platform for Large-Scale Grids,''
\emph{IEEE Transactions on Smart Grid}, vol.~8, no.~3, pp.~1417--1426, May 2017.
doi: \href{https://doi.org/10.1109/TSG.2016.2606888} {10.1109/TSG.2016.2606888}.

\bibitem{veragrid}
S. Peñate Vera, J. Fanals i Batllori, C. Alegre, benceszirbik, A. Blanco Castro, 
M. Lavoie, Q. Moya, L. Raiyan, B. Lüers, fernpos, M. Zsebeházi, J. Soler, 
ManuelNvro, ramferan, C. E. Fray, jag0nzalez, J. Gorcs, mmutto, 
peterkulik-navitasoft, R. Yu, 
``SanPen/VeraGrid: 5.4.0 VeraGrid first release (v5.4.0-veragrid),'' Zenodo, 2025. 
DOI: 10.5281/zenodo.16971840. [Online]. Available: \url{https://doi.org/10.5281/zenodo.16971840}

\bibitem{pena2017extended}
I.~Pena, C.~B.~Martinez-Anido, and B.-M.~Hodge,
``An extended IEEE 118-bus test system with high renewable penetration,''
\emph{IEEE Transactions on Power Systems}, vol.~33, no.~1, pp.~281--289, 2018
doi: \href{https://doi.org/10.1109/TPWRS.2017.2695963}{10.1109/TPWRS.2017.2695963}.

\bibitem{pycompss_softx}
R.~M.~Badia, J.~Conejero, C.~Díaz, J.~Ejarque, D.~Lezzi, F.~Lordan,
C.~Ramón-Cortés, and R.~Sirvent,
``COMP Superscalar, an interoperable programming framework,''
\emph{SoftwareX}, vol.~3--4, pp.~32--36, Dec.~2015.
doi: \href{https://doi.org/10.1016/j.softx.2015.10.004} {10.1016/j.softx.2015.10.004}.

\bibitem{pycompss_tejedor2017}
E. Tejedor, Y. Becerra, G. Alomar, A. Queralt, R. M. Badia, J. Torres, T. Cortes, and J. Labarta,
“PyCOMPSs: Parallel Computational Workflows in Python,”
\textit{The International Journal of High Performance Computing Applications},
vol. 31, no. 1, pp. 66--82,
Jan. 2017,
doi: \href{https://doi.org/10.1177/1094342015594678}{10.1177/1094342015594678}.

\bibitem{pal2006}
B. C. Pal and B. Chaudhuri, \emph{Robust Control in Power Systems}, Springer, 2005.
doi: \href{https://doi.org/10.1007/b136490}{10.1007/b136490}.

\bibitem{sperstad_vulnerability_graph}
I. B. Sperstad, E. H. Solvang and S. H. Jakobsen,
“A graph-based modelling framework for vulnerability analysis of critical sequences of events in power systems,”
\emph{International Journal of Electrical Power \& Energy Systems}, 
vol.~124, art.~106408, Feb. 2021,
doi: \href{https://doi.org/10.1016/j.ijepes.2020.106408}{10.1016/j.ijepes.2020.106408}.

\bibitem{tarjan}
R.~Tarjan,
\emph{Depth-first search and linear graph algorithms},
SIAM Journal on Computing, \textbf{1}(2), pp. 146--160, 1972.
doi: \href{https://doi.org/10.1137/0201010}{10.1137/0201010}.

\bibitem{nord4_bsc}
Barcelona Supercomputing Center (BSC),
``Nord4 overview,''
[Online]. Available: \url{https://www.bsc.es/supportkc/docs/Nord4/overview/}
[Accessed: Nov.~20, 2025].




\end{thebibliography}
\vfill

\end{document}
